\documentclass[11pt]{article}

\usepackage[margin=1in]{geometry}
\usepackage[T1]{fontenc}
\usepackage{lmodern}
\usepackage{microtype}
\usepackage{amsmath,amssymb}
\usepackage{booktabs}
\usepackage{tabularx}
\usepackage{longtable}
\usepackage{array}
\usepackage{enumitem}
\usepackage{xcolor}
\usepackage{hyperref}

\hypersetup{
  colorlinks=true,
  linkcolor=blue!50!black,
  urlcolor=blue!50!black,
  citecolor=blue!50!black
}

\setlist[itemize]{leftmargin=1.4em}
\setlist[enumerate]{leftmargin=1.6em}
\renewcommand{\arraystretch}{1.15}

\title{\textbf{A Cross-OSDR Network-Contrast Framework for Kidney Spaceflight Aging}\\
\large Working methodology memo}
\author{RRRM-2 Kidney Transcriptome Project}
\date{May 2026}

\begin{document}
\maketitle

\begin{abstract}
This memo translates the current project debate into a concrete analysis plan.
The goal is not to replace the original network-rewiring question with a generic
WGCNA story. The goal is to keep the original estimand---whether spaceflight
amplifies, dampens, reverses, or redirects age-associated kidney network
change---while moving from a single underpowered dataset to a cross-OSDR
network-contrast framework. RRRM-2 remains the only exact Age $\times$ Arm
$\times$ Flight kidney RNA-seq design. Other OSDR kidney datasets cannot
replicate the full factorial question, but they can test whether RRRM-2 flight
network directions, node-level scores, module-level programs, and silent-shifter
candidates recur across independent missions, strains, sexes, durations, and
platforms.
\end{abstract}

\tableofcontents
\newpage

\section{Plain-English Summary}

The original project should not be thrown away. The original idea is still the
right scientific question:

\begin{quote}
Does spaceflight alter age-related kidney transcriptional network organization,
and does it do so differently in ISS-T and LAR?
\end{quote}

The problem is that RRRM-2 alone has $n=5$ per Age $\times$ Arm $\times$ Flight
cell. That is enough to estimate some module- or pathway-scale effects, but it is
not enough to make individual edge discovery carry the paper. Therefore the
framework should shift from:

\begin{quote}
\emph{single-dataset edge discovery}
\end{quote}

to:

\begin{quote}
\emph{multi-cohort network contrast estimation, projection, and replication at
edge, gene, module, and pathway resolution.}
\end{quote}

The key idea is simple: a contrast can be represented as a vector. The vector can
be an expression vector, but for this project it should primarily be a
\emph{network} vector:

\begin{itemize}
  \item an edge-effect vector;
  \item a gene differential-connectivity vector;
  \item a centrality-change vector;
  \item a delta-kME vector;
  \item a LIONESS-derived node-score vector;
  \item or a module/pathway activity vector.
\end{itemize}

Then the question ``amplifies, dampens, or redirects'' becomes a geometry
question. If the flight-aging network vector points in the same direction as the
control-aging network vector, flight is aging-like or amplifying. If it points in
the opposite direction, flight is reversing. If most of it is orthogonal, flight
is redirecting the kidney into a different network state.

\section{What This Is Not}

This plan is not:

\begin{itemize}
  \item a generic WGCNA paper;
  \item a DE-only or pathway-only paper;
  \item a claim that Tabula Muris Senis droplet data should be used to build raw
        co-expression networks;
  \item a raw pooled mega-analysis that ignores mission, strain, sex, platform,
        duration, or dissection protocol;
  \item another attempt to force $n=5$ separate condition-level correlation
        networks to produce stable individual edge discoveries.
\end{itemize}

This plan is:

\begin{itemize}
  \item a way to preserve the original network-rewiring question;
  \item a way to use RRRM-2 as the only exact factorial bridge;
  \item a way to use other OSDR kidney datasets as independent flight-network
        replication or context cohorts;
  \item a way to decide which claims survive beyond RRRM-2.
\end{itemize}

\section{How We Got Here}

This section is included because the current proposal can sound, at first, like
a sudden pivot into dot products or embeddings. It is not. It is the result of
walking the original question through several methods and keeping only the parts
that still answer the biology honestly.

\subsection{Starting point: age-dependent network rewiring}

The original project was not ``do differential expression and then WGCNA.'' The
original project was:

\begin{quote}
Use the RRRM-2 factorial design to ask whether spaceflight changes the way mouse
kidney transcriptional networks age, separately for ISS-T and LAR, after
accounting for kidney cell composition.
\end{quote}

That implies several desired outputs:

\begin{itemize}
  \item estimate age-associated network change in controls;
  \item estimate age-associated network change under flight;
  \item test whether flight amplifies, dampens, reverses, or redirects the aging
        network direction;
  \item do this separately for ISS-T and LAR;
  \item find genes with high network change but modest differential expression;
  \item interpret those genes with targeted kidney biology, not generic pathway
        lists;
  \item validate signals with sample-specific networks, external cohorts, or
        metadata-aware projections.
\end{itemize}

In statistical terms, the original estimand was always a contrast of contrasts:

\[
  \left(\mathrm{Old}_{FLT}-\mathrm{Young}_{FLT}\right)
  -
  \left(\mathrm{Old}_{GC}-\mathrm{Young}_{GC}\right).
\]

For this project, the object inside each parenthesis is not only expression. It
can be a network edge, node score, module score, pathway score, or
sample-specific network summary.

\subsection{First constraint: separate condition networks are not credible}

The most direct version of the idea would build one network for each condition:
Young GC, Old GC, Young FLT, Old FLT, separately inside ISS-T and LAR. That would
match the biology neatly, but RRRM-2 has only about $n=5$ per cell. At that size,
condition-specific correlations are too unstable to be the main inferential
object. This forced the first methodological move:

\begin{quote}
Do not estimate completely separate networks per condition. Estimate effects on
a shared or constrained network coordinate system.
\end{quote}

That is why shared-topology networks, fixed edge universes, shrinkage
correlations, differential connectivity summaries, and LIONESS-style
sample-specific networks entered the project.

\subsection{Second move: WGCNA as a stability and annotation layer}

WGCNA entered for a sensible reason. If individual edges are unstable, modules
can provide a more stable biological resolution. WGCNA can define co-expression
modules, eigengenes, hub genes, and module-trait associations. That is useful.

The problem is that WGCNA does not naturally answer the original factorial
network question. Its module eigengene is the first principal component of a
module's expression matrix. That eigengene can detect a Flight $\times$ Age or
Flight $\times$ Arm pattern if the relevant signal lies on the dominant module
axis. But if the interaction lives in a smaller-variance within-module direction,
the eigengene can blur it away. Also, WGCNA module-trait association is not the
same as edge-level or node-level rewiring.

Therefore WGCNA was demoted from:

\begin{quote}
primary proof of age-dependent rewiring
\end{quote}

to:

\begin{quote}
secondary structure, compression, and biological annotation.
\end{quote}

That demotion is not anti-WGCNA. It is matching the method to the question.

\subsection{Third move: LIONESS and sample-specific networks}

LIONESS entered because the project needed sample-level network summaries. A
sample-specific network gives each mouse a network feature vector, which can then
be modeled with ordinary factorial machinery:

\[
  \mathrm{network\ feature} \sim
  \mathrm{Age} \times \mathrm{Flight} \times \mathrm{Arm}
  + \mathrm{covariates}.
\]

This is closer to the original design than a single pooled WGCNA eigengene test.
It lets the model ask whether a network feature changes with age differently
under flight and whether that differs between ISS-T and LAR.

The cost is multiple testing. Edge-wise LIONESS models across many edges are
underpowered at RRRM-2 sample size, so strict FDR can be sparse or null. That
does not make the method wrong. It means individual edge discovery cannot carry
the manuscript alone.

\subsection{Fourth move: silent shifters}

The silent-shifter idea was created to capture a biologically plausible pattern:

\begin{quote}
Some genes may not change much in mean expression, but their network
neighborhood, connectivity, or regulatory context may change substantially.
\end{quote}

That is still a good idea. The issue is that a single-study silent-shifter list
can be fragile. A stronger version treats silent shifting as a ranking or
recurrence phenomenon:

\begin{quote}
genes repeatedly high in network-change rank and low in differential-expression
rank across strata or studies.
\end{quote}

So silent shifters remain useful, but as reproducible node-level candidates and
pathway/module annotations rather than as a hard single-study gene list.

\subsection{Fifth move: node2vec and latent topology}

Node2vec entered for the same reason silent shifters entered: it seemed like a
way to measure whether a gene's network neighborhood moved even when expression
did not. In principle, if a gene shifts from one network neighborhood to another,
its learned graph embedding can move.

The problem is interpretability. A node2vec coordinate is a latent coordinate.
Its axes are not fixed biological quantities, embeddings from different graphs
must be aligned, and results can depend on graph construction, random-walk
parameters, embedding dimension, and random seed. That makes node2vec reasonable
as an exploratory description of topology movement, but weak as the primary
factorial estimand.

This is why the current plan avoids saying ``make embeddings and compare them.''
The preferred language is:

\begin{quote}
estimate interpretable network contrast vectors, then decompose their direction,
magnitude, and orthogonal redirected component.
\end{quote}

Node2vec can still be used as a sensitivity or exploratory layer, but not as the
main claim.

\subsection{Sixth move: the grey60/WGCNA convergence attempt}

One attempted rescue was to make WGCNA grey60 the biological center of the paper
and ask whether LIONESS-factorial genes converged into that module. The
repository results did not support that claim. The 11 ISS-T-Young LIONESS
candidate genes did not concentrate in grey60. The best overlap was marginal in
purple, and grey60 had zero overlap with those candidates.

At the eigengene level, grey60 did show an ISS-T-specific pattern, but so did
royalblue, salmon, and pink. Therefore the correct inference is not:

\begin{quote}
one unique grey60 ECM module drives the entire response.
\end{quote}

The more honest inference is:

\begin{quote}
RRRM-2 contains an ISS-T-biased module-level signal across several biological
programs, while gene-level network rankings are method-dependent.
\end{quote}

That result was important because it prevented the manuscript from becoming a
WGCNA-first story that no longer matched the original question.

\subsection{Seventh move: OSD-253 as a proposed TLR4 discriminator}

OSD-253 was considered as a possible validation or mechanism test because it
contains C57BL/6J and C3H/HeJ mice, and C3H/HeJ carries impaired TLR4 signaling.
If grey60 were uniquely TLR4-dependent, OSD-253 might have shown selective
attenuation of grey60-like module projection in C3H/HeJ.

The metadata made this more complicated. OSD-253 has original controls and rerun
controls with light wavelength, read length, and run-context differences.
Therefore it cannot be used casually as a clean TLR4 validation. A
metadata-aware projection is useful, but only as scenario-based context. The
current repository analysis did not provide a clean grey60-specific C3H
attenuation result.

This means OSD-253 should not decide the whole paper. It should constrain and
contextualize mechanistic claims.

\subsection{Current synthesis: contrast-vector decomposition}

The current proposal is the cleaned-up version of the original methodology. It
keeps the original biological question but changes the level at which the
statistics are expected to be reliable.

\begin{center}
\begin{tabularx}{0.92\textwidth}{p{0.30\textwidth}X}
\toprule
Earlier version & Current version \\
\midrule
Find individual rewired edges in RRRM-2. &
Estimate global, module, pathway, and node-level network contrast vectors first;
treat individual edges as exploratory. \\
\addlinespace
Use WGCNA as primary evidence. &
Use WGCNA as module definition, compression, and annotation. \\
\addlinespace
Use node2vec as central topology movement metric. &
Use node2vec only as exploratory; prefer interpretable contrast vectors. \\
\addlinespace
Validate one grey60 mechanism in OSD-253. &
Use OSD-253 as metadata-aware strain/duration context, not a clean decisive
validation. \\
\addlinespace
RRRM-2 must succeed alone. &
RRRM-2 defines the exact factorial contrast; external OSDR cohorts test whether
flight-network directions recur. \\
\bottomrule
\end{tabularx}
\end{center}

So the current framework is not a retreat from the original project. It is the
same project expressed at a statistically defensible resolution.

\section{Scientific Question}

\subsection{Original question}

The original question can be written as:

\begin{quote}
After accounting for kidney cell composition, does spaceflight amplify, dampen,
reverse, or redirect age-associated kidney transcriptional network change,
separately in ISS-T and LAR?
\end{quote}

The key feature is that this is a network contrast question, not simply a module
eigengene association question.

\subsection{What WGCNA can and cannot answer}

WGCNA is useful for defining stable co-expression modules, module eigengenes, hub
genes, and module-trait associations. It can support statements such as:

\begin{quote}
This module is flight-associated, age-associated, or enriched for DCT, ECM, ER
stress, transport, immune, or mitochondrial biology.
\end{quote}

But WGCNA is not the native method for asking whether the age-related
\emph{network contrast} changes under flight. Standard module eigengene models can
miss genes whose average expression is stable but whose network neighborhood
changes. Therefore WGCNA should be a secondary annotation and compression layer,
not the primary estimand.

\section{Dataset Roles}

\begin{longtable}{p{0.14\textwidth}p{0.20\textwidth}p{0.36\textwidth}p{0.20\textwidth}}
\toprule
Dataset & Assay / tissue & Design role & Use in this framework \\
\midrule
\endfirsthead
\toprule
Dataset & Assay / tissue & Design role & Use in this framework \\
\midrule
\endhead
OSD-771 / GLDS-674 &
Kidney bulk RNA-seq &
RRRM-2. Female C57BL/6NTac. Spaceflight, age, euthanasia location / arm.
Young and old; ISS-T and LAR; FLT, GC, VIV, BSL. &
Primary exact Age $\times$ Arm $\times$ Flight network-contrast dataset. \\
\addlinespace
OSD-513 / GLDS-513 &
Kidney bulk RNA-seq &
RR-23. Male C57BL/6J. 20 FLT, 20 HGC, 20 VGC; no age factor. &
Best powered flight network-response anchor. Validate whether RRRM-2 flight
network directions recur in a larger cohort. \\
\addlinespace
OSD-253 / GLDS-253 &
Kidney bulk RNA-seq &
RR-7. Female C57BL/6J and C3H/HeJ. 25 and 75 day durations. Original GC and GC
rerun controls have important light/read-length complications. &
Strain and duration context. Do not treat as a clean TLR4 validation without
metadata-aware sensitivity analyses. \\
\addlinespace
OSD-102 / GLDS-102 &
Kidney bulk RNA-seq plus multi-omics &
RR-1. Female C57BL/6J, 16 weeks, FLT versus GC. No age factor. &
Flight network context and replication of module/pathway/node directions. \\
\addlinespace
OSD-163 / GLDS-163 &
Kidney bulk RNA-seq plus multi-omics &
RR-3. Female BALB/c, 12 weeks, FLT versus GC. No age factor. &
Strain/context replication. Useful for whether flight network programs are
C57-specific or broader. \\
\addlinespace
OSD-462 / GLDS-462 &
Kidney transcriptomics, proteomics, phosphoproteomics &
RR-10. Female WT / p21-null design; WT subset can be used carefully. No age
factor. &
Mechanistic context for stress, cell-cycle, p21, and proteomic support. \\
\addlinespace
OSD-457 / GLDS-457 &
Multi-tissue RNA-seq including kidney &
JAXA MHU-3. WT and Nrf2-KO, spaceflight and genotype. No age factor. &
Stress-response and genotype context, especially Nrf2/oxidative stress. \\
\addlinespace
OSD-913 / GLDS-735 &
Kidney miRNA-seq &
RRRM-1 / RR-8. Age, flight, euthanasia location, dissection condition. Not
bulk mRNA. &
Regulatory-layer age $\times$ flight context. Do not use as mRNA co-expression
network evidence. \\
\addlinespace
Tabula Muris Senis / aging references &
Single-cell or bulk aging atlases &
Independent aging references. &
Use for deconvolution, cell-type aging signatures, and annotation. Do not build
raw droplet co-expression networks as the primary network reference. \\
\bottomrule
\end{longtable}

\section{Core Notation}

Let $X_s$ be the residualized expression matrix for study $s$, with genes
$g \in G$ and samples $i \in I_s$.

\subsection{Residualized expression}

Within each study, residualize expression using metadata-appropriate covariates:

\[
X_{g i} = \mu_g + C_i \alpha_g + T_i \gamma_g + \varepsilon_{g i}
\]

where $C_i$ includes estimated kidney cell-type composition and $T_i$ includes
technical covariates such as batch, library batch, RIN/RINe, read length,
sequencing run, and study-specific confounds. The residuals
$\hat{\varepsilon}_{g i}$ are the input for network estimation.

\subsection{Network vector}

Define a frozen edge universe $E = \{(g,h)\}$ on a common gene panel. For a
contrast $c$, define an edge-effect vector:

\[
\mathbf{e}_{s,c} =
\left( \Delta w_{s,c,1}, \Delta w_{s,c,2}, \ldots, \Delta w_{s,c,|E|} \right)
\]

where $\Delta w_{s,c,k}$ is the contrast effect for edge $k$ in study $s$.

Also define gene-level node vectors:

\[
\mathbf{n}_{s,c} =
\left( \Delta N_{s,c,1}, \Delta N_{s,c,2}, \ldots, \Delta N_{s,c,|G|} \right)
\]

where $\Delta N$ can be differential connectivity, centrality change, delta-kME,
or an aggregated LIONESS node score.

\section{Backbone Strategy}

\subsection{Why a shared backbone is needed}

Separate $n=5$ correlation networks are unstable. A raw Pearson correlation has
Fisher-$z$ standard error approximately

\[
\mathrm{SE}(z) = \frac{1}{\sqrt{n-3}}.
\]

At $n=5$, this is about $0.71$, before even taking differences between
conditions. Therefore condition-specific networks cannot be the primary object.

\subsection{Frozen candidate edge universe}

Use a frozen candidate edge universe rather than testing all gene pairs. Options:

\begin{enumerate}
  \item \textbf{Control-only shrinkage backbone:} build a Ledoit-Wolf / shrinkage
  partial-correlation backbone from pooled residualized controls, excluding FLT
  labels from edge selection.
  \item \textbf{Multi-study control backbone:} build study-specific control
  backbones, then take the union or recurring top-$k$ edges.
  \item \textbf{Biology-constrained backbone:} restrict to edges within or
  between kidney-relevant pathways/modules, such as DCT/transport, ECM, ER
  stress, mitochondria, immune, lipid metabolism, and tubular injury.
\end{enumerate}

The backbone is not itself the inferential result. It is the coordinate system
for contrast estimation.

\section{RRRM-2: Exact Primary Estimands}

RRRM-2 is the only exact dataset for the original question. Define network
vectors separately by arm.

\subsection{Control aging network vector}

For arm $a \in \{\mathrm{ISS\mbox{-}T}, \mathrm{LAR}\}$:

\[
\mathbf{A}^{a}_{GC}
= \mathbf{N}(\mathrm{Old}, a, GC)
- \mathbf{N}(\mathrm{Young}, a, GC)
\]

This is the age-associated network change in controls.

\subsection{Flight aging network vector}

\[
\mathbf{A}^{a}_{FLT}
= \mathbf{N}(\mathrm{Old}, a, FLT)
- \mathbf{N}(\mathrm{Young}, a, FLT)
\]

This is the age-associated network difference under flight.

\subsection{Spaceflight modification of age rewiring}

\[
\Delta^{a}_{AgeRewire}
= \mathbf{A}^{a}_{FLT} - \mathbf{A}^{a}_{GC}
\]

This is the within-arm Flight $\times$ Age network interaction. The arm
difference is:

\[
\Delta_{Arm}
= \Delta^{ISS\mbox{-}T}_{AgeRewire}
- \Delta^{LAR}_{AgeRewire}.
\]

\section{Amplify, Dampen, Reverse, Redirect}

This is where the dot product enters. The dot product is not replacing network
analysis; it is a way to classify the direction of one network contrast relative
to another.

\subsection{Projection coefficient}

For a given arm:

\[
\beta_a =
\frac{\mathbf{A}^{a}_{FLT} \cdot \mathbf{A}^{a}_{GC}}
     {\mathbf{A}^{a}_{GC} \cdot \mathbf{A}^{a}_{GC}}.
\]

Interpretation:

\begin{itemize}
  \item $\beta_a > 1$: flight amplifies the control aging network direction.
  \item $0 < \beta_a < 1$: flight dampens the control aging network direction.
  \item $\beta_a < 0$: flight reverses the control aging network direction.
  \item $\beta_a \approx 0$: flight is mostly not aligned with the control aging
        network direction.
\end{itemize}

\subsection{Cosine similarity}

\[
\cos(\theta_a) =
\frac{\mathbf{A}^{a}_{FLT} \cdot \mathbf{A}^{a}_{GC}}
     {\|\mathbf{A}^{a}_{FLT}\| \, \|\mathbf{A}^{a}_{GC}\|}.
\]

This measures directional similarity independent of magnitude.

\subsection{Orthogonal / redirected component}

\[
\mathbf{R}_a =
\mathbf{A}^{a}_{FLT} - \beta_a \mathbf{A}^{a}_{GC}.
\]

\[
\rho_a =
\frac{\|\mathbf{R}_a\|}{\|\mathbf{A}^{a}_{FLT}\|}.
\]

Interpretation:

\begin{itemize}
  \item small $\rho_a$: flight aging pattern is mostly along the control aging
        axis;
  \item large $\rho_a$: flight aging pattern is redirected into a different
        network direction.
\end{itemize}

\section{Inference}

\subsection{Permutation tests}

For RRRM-2, use label permutations that respect the design:

\begin{itemize}
  \item permute FLT/GC labels within Age $\times$ Arm when testing flight
        effects;
  \item permute Young/Old labels within Arm $\times$ Environment when testing
        age effects;
  \item preserve cage/batch constraints where metadata requires it;
  \item recompute contrast vectors, projections, and orthogonal fractions under
        each permutation.
\end{itemize}

Primary permutation statistics:

\begin{itemize}
  \item $\beta_a$ for amplification/dampening;
  \item $\cos(\theta_a)$ for directional alignment;
  \item $\rho_a$ for redirectedness;
  \item $\|\Delta^{a}_{AgeRewire}\|$ for interaction magnitude;
  \item module/pathway concentration of $\mathbf{R}_a$.
\end{itemize}

\subsection{Avoiding overclaiming}

Individual edge FDR should not be the sole success criterion. At RRRM-2 cell
sizes, edge-wise FDR is expected to be sparse or null. The hierarchy should be:

\begin{enumerate}
  \item global network-vector statistic;
  \item module/pathway concentration;
  \item gene/node ranking;
  \item individual edges only as exploratory annotations.
\end{enumerate}

\section{Cross-OSDR Extension}

Other OSDR studies do not generally have age $\times$ arm structure. Therefore
they cannot replicate the full RRRM-2 estimand. They can replicate pieces.

\subsection{Within-study flight network vectors}

For each external study $s$, define:

\[
\mathbf{F}_s =
\mathbf{N}_s(FLT) - \mathbf{N}_s(GC/HGC).
\]

Where possible, define study-specific variants:

\begin{itemize}
  \item OSD-253: $\mathbf{F}_{strain,duration,control}$ for strain, duration,
        and original versus rerun controls;
  \item OSD-513: powered male C57 flight vector;
  \item OSD-102: female C57 flight vector;
  \item OSD-163: BALB/c flight vector;
  \item OSD-462: WT subset flight vector;
  \item OSD-457: WT and Nrf2-KO flight vectors if sample sizes permit.
\end{itemize}

\subsection{Cross-study projection tests}

Test whether external flight network vectors align with RRRM-2-derived vectors:

\[
\cos(\mathbf{F}_s, \mathbf{F}_{RRRM2,ISS})
\]

\[
\cos(\mathbf{F}_s, \Delta^{ISS\mbox{-}T}_{AgeRewire})
\]

\[
\cos(\mathbf{F}_s, \mathbf{R}_{ISS})
\]

This asks whether the RRRM-2 network direction is recurring across missions, not
whether every external dataset has the same factorial design.

\subsection{Meta-analysis}

For gene/node-level scores, combine study-level signed statistics using:

\begin{itemize}
  \item signed Stouffer $Z$ with weights proportional to effective sample size;
  \item random-effects meta-analysis for module/pathway scores;
  \item rank aggregation for silent-shifter candidates;
  \item leave-one-study-out robustness checks.
\end{itemize}

Do not pool raw expression across missions as the primary analysis. Pool
contrast-level statistics.

\section{Silent Shifters}

The single-study silent-shifter definition is fragile. In the cross-OSDR
framework, silent shifters should become a reproducibility concept.

\subsection{Within-study score}

For each gene $g$:

\[
S_{g,s} = \mathrm{network\_change}_{g,s}
\]

\[
D_{g,s} = |\mathrm{DE}_{g,s}|
\]

A gene is a candidate silent shifter in study $s$ if it has high $S_{g,s}$ and
low $D_{g,s}$.

\subsection{Cross-study criterion}

A stronger definition:

\begin{quote}
A reproducible silent shifter is a gene with consistently high network-change
rank across studies or strata, while remaining modest in differential-expression
rank.
\end{quote}

This should be tested with enrichment and rank aggregation, not treated as a
single-study gene list with strong individual-gene claims.

\section{Role of WGCNA and Modules}

WGCNA should remain in the project, but its role should be explicit:

\begin{itemize}
  \item define or annotate stable residual-expression modules;
  \item compress network and expression signals into module-level summaries;
  \item ask whether high network-vector loadings concentrate in modules;
  \item interpret redirected components biologically;
  \item test whether RRRM-2 module signals project into external cohorts.
\end{itemize}

WGCNA should not be asked to prove age-dependent topology rewiring.

\section{Tabula Muris Senis and Aging References}

Tabula Muris Senis should not be used as a raw droplet co-expression network
source for this paper. It can be used in three safer ways:

\begin{enumerate}
  \item \textbf{Cell-type deconvolution reference:} the current use case.
  \item \textbf{Cell-type-specific aging annotations:} pseudobulk within cell
        type and age where metadata supports animal-level aggregation.
  \item \textbf{Biological labels:} identify whether network-loaded genes are PT,
        DCT, LOH, collecting duct, endothelial, immune, fibroblast, or other
        cell-type-biased aging genes.
\end{enumerate}

If a good independent bulk kidney aging reference is available, use it to provide
an external aging axis. But that should be secondary to the network aging axis
estimated inside RRRM-2.

\section{Metadata-Specific Rules}

\subsection{OSD-253}

OSD-253 requires separate scenarios:

\begin{itemize}
  \item original GC: matched read length, but FLT is white-light and original GC
        is blue-light;
  \item white-light rerun GC: fixes light wavelength, but changes read length and
        run context;
  \item blue-light rerun GC: sensitivity only.
\end{itemize}

Any strain/TLR4 statement must survive metadata-aware scenario analysis.

\subsection{OSD-771}

OSD-771 is the primary factorial bridge. Euthanasia location / arm must be treated
as part of the biological and procedural design:

\begin{itemize}
  \item ISS-T FLT: terminal on-orbit collection;
  \item LAR FLT: return to Earth and recovery;
  \item controls are Earth-collected;
  \item therefore ISS-T-specific signals cannot be casually called purely
        ``acute flight'' without discussing collection protocol.
\end{itemize}

\section{Proposed Analysis Phases}

\subsection*{Phase 0: Dataset inventory}

\begin{itemize}
  \item Download or verify VST/count matrices and ISA metadata for OSD-771,
        OSD-513, OSD-253, OSD-102, OSD-163, OSD-462, OSD-457, and OSD-913.
  \item Create a harmonized metadata table with study, strain, sex, age, duration,
        FLT/GC/VIV/BSL, arm/euthanasia, platform, read length, library batch, RIN,
        and control scenario.
\end{itemize}

\subsection*{Phase 1: Harmonization and residualization}

\begin{itemize}
  \item Map genes to common Ensembl IDs.
  \item Compute cell-type deconvolution using the existing kidney reference.
  \item Fit study-specific residualization models.
  \item Save residual matrices with provenance.
\end{itemize}

\subsection*{Phase 2: Frozen backbone}

\begin{itemize}
  \item Define one or more frozen edge universes.
  \item Record exactly how edges were selected.
  \item Lock the backbone before contrast interpretation.
\end{itemize}

\subsection*{Phase 3: Within-study network contrasts}

\begin{itemize}
  \item RRRM-2: estimate full Age $\times$ Arm $\times$ Flight network contrasts.
  \item OSD-513 and other cohorts: estimate FLT versus GC/HGC network vectors.
  \item OSD-253: estimate strain-, duration-, and control-scenario-specific
        network vectors.
\end{itemize}

\subsection*{Phase 4: Projection and redirection}

\begin{itemize}
  \item Compute $\beta$, cosine similarity, and orthogonal fraction for RRRM-2.
  \item Compute cross-study alignment of external flight vectors with RRRM-2
        network vectors.
  \item Use permutation or bootstrap for uncertainty.
\end{itemize}

\subsection*{Phase 5: Gene, module, and pathway interpretation}

\begin{itemize}
  \item Aggregate edges to genes.
  \item Identify reproducible node-level contributors.
  \item Test module/pathway enrichment of aligned and redirected components.
  \item Re-evaluate silent shifters as recurring rank patterns.
\end{itemize}

\subsection*{Phase 6: Decision}

\begin{longtable}{p{0.36\textwidth}p{0.56\textwidth}}
\toprule
Observed result & Manuscript implication \\
\midrule
\endfirsthead
\toprule
Observed result & Manuscript implication \\
\midrule
\endhead
RRRM-2 age-rewiring vector is significant and recurs across external flight
vectors. &
Strong network-aging paper. Spaceflight modifies kidney aging network
architecture. \\
\addlinespace
RRRM-2 is weak alone, but external flight vectors align with RRRM-2 flight or
redirected network programs. &
Cross-OSDR flight-network biology paper. RRRM-2 provides the factorial bridge;
external cohorts provide recurrence. \\
\addlinespace
Only WGCNA module activity is positive; network vectors do not recur. &
Do not force a network-aging paper. Consider a modest module-activity paper or
stop. \\
\addlinespace
No coherent module, pathway, node, or projection signal survives. &
Do not write RRRM-2 as a standalone biological network paper. Use as preliminary
or negative constraint evidence. \\
\bottomrule
\end{longtable}

\section{Minimum Viable Version}

If time is limited, the minimum viable version is:

\begin{enumerate}
  \item RRRM-2 exact network-vector redirection test using the existing
        residualized/LIONESS/direct-correlation machinery.
  \item OSD-513 powered FLT versus HGC network vector.
  \item OSD-253 metadata-aware strain/duration sensitivity.
  \item Compare vector alignment and module/pathway concentration across these
        three.
\end{enumerate}

This is enough to determine whether the cross-OSDR network paper is alive.

\section{Bottom Line}

The original methodology is not being discarded. It is being reframed from:

\begin{quote}
``RRRM-2 alone discovers age-dependent network rewiring genes.''
\end{quote}

to:

\begin{quote}
``RRRM-2 defines the only exact age-by-arm flight network contrast, and a
cross-OSDR framework tests whether that network direction recurs across
independent kidney spaceflight cohorts.''
\end{quote}

That is closer to the original ambition than a WGCNA-only paper, and it avoids
pretending that $n=5$ per factorial cell can support strong individual edge
discovery by itself.

\section*{Useful Source Links}

\begin{itemize}
  \item OSD-771 / RRRM-2 kidney RNA-seq:
  \url{https://osdr.nasa.gov/bio/repo/data/studies/OSD-771}
  \item OSD-513 / RR-23 kidney RNA-seq:
  \url{https://osdr.nasa.gov/bio/repo/data/studies/OSD-513}
  \item OSD-253 / RR-7 kidney RNA-seq:
  \url{https://osdr.nasa.gov/bio/repo/data/studies/OSD-253}
  \item OSD-102 / RR-1 kidney:
  \url{https://osdr.nasa.gov/bio/repo/data/studies/OSD-102}
  \item OSD-163 / RR-3 kidney:
  \url{https://osdr.nasa.gov/bio/repo/data/studies/OSD-163}
  \item OSD-462 / RR-10 kidney:
  \url{https://osdr.nasa.gov/bio/repo/data/studies/OSD-462}
  \item OSD-457 / MHU-3 multi-tissue:
  \url{https://osdr.nasa.gov/bio/repo/data/studies/OSD-457}
  \item OSD-913 / RRRM-1 kidney miRNA:
  \url{https://osdr.nasa.gov/bio/repo/data/studies/OSD-913}
\end{itemize}

\end{document}
